\section{Fundamentos y estado del arte}

\paragraph{AE.} Minimiza el error de reconstrucción
$\mathcal{L}_{AE}=\lVert x-g_\theta(f_\phi(x))\rVert_1$. Su ventaja es la
simplicidad; su debilidad es que el cuello de botella no impone una distribución
latente y puede generalizar la reconstrucción fuera de la normalidad.

\paragraph{VAE.} Aproxima $q_\phi(z\mid x)$ mediante una normal parametrizada y
optimiza
\[
\mathcal{L}_{VAE}=\lVert x-\hat{x}\rVert_1+
\beta D_{KL}\!\left(q_\phi(z\mid x)\,\|\,\mathcal{N}(0,I)\right).
\]
El truco de reparametrización permite derivar a través del muestreo
\cite{kingma2014vae}. Estudios recientes comparan variantes convolucionales y
transformers y confirman que la arquitectura y $\beta$ condicionan el resultado
\cite{nguyen2024vae}.

\paragraph{GANomaly.} Frente a la optimización iterativa de AnoGAN, GANomaly
aprende dos codificadores y un decodificador. Su generador produce
$z=E_1(x)$, $\hat{x}=G(z)$ y $\hat{z}=E_2(\hat{x})$. Combina pérdida contextual,
latente y adversarial; durante inferencia, $\lVert z-\hat{z}\rVert_1$ actúa como
puntuación \cite{akcay2019ganomaly}. La comparación es relevante porque el
entrenamiento adversarial puede mejorar representaciones, pero añade
inestabilidad y coste.

Los benchmarks industriales como MVTec AD impulsaron protocolos comunes
\cite{bergmann2019mvtec}. En medicina, BMAD extiende esa idea a seis conjuntos
de cinco dominios y quince algoritmos \cite{bao2024bmad}. Este trabajo no busca
reproducir todo el leaderboard: aísla las tres familias solicitadas en la
propuesta y controla las variables experimentales.
